\documentclass[11pt]{article}
\usepackage[a4paper,margin=18mm]{geometry}
\usepackage{fontspec}
\setmainfont{Latin Modern Roman}
\usepackage{amsmath,amssymb,amsthm,amscd,graphicx,color}
\usepackage{xurl}
\usepackage[unicode,hidelinks]{hyperref}
\allowdisplaybreaks
\setlength\emergencystretch{3em}
\numberwithin{equation}{section}


\newtheorem{Theorem}{Theorem}[section]
\newtheorem{Corollary}[Theorem]{Corollary}
\newtheorem{Lemma}[Theorem]{Lemma}
 { \theoremstyle{definition}
\newtheorem{Remark}[Theorem]{Remark} }

\def\Xint#1{\mathchoice
{\XXint\displaystyle\textstyle{#1}}%
{\XXint\textstyle\scriptstyle{#1}}%
{\XXint\scriptstyle\scriptscriptstyle{#1}}%
{\XXint\scriptscriptstyle\scriptscriptstyle{#1}}%
\!\int}
\def\XXint#1#2#3{{\setbox0=\hbox{$#1{#2#3}{\int}$}
\vcenter{\hbox{$#2#3$}}\kern-.5\wd0}}
\def\ddashint{\Xint=}
\def\dashint{\Xint-}


\newcommand{\intZ}{\mathbb{Z}}
\newcommand{\prodprime}{\sideset{}{'}\prod}
\newcommand{\lHopital}{l'H\^{o}pital}
\newcommand{\Weylchamber}{\mathbb{W}}
\newcommand{\bigO}{\mathcal{O}}

\DeclareMathOperator*{\Res}{Res}
\DeclareMathOperator{\id}{id}


\begin{document}
\begin{center}\Large Finite-particle q-boson transition probabilities with positive site-dependent conductances on the integer lattice\end{center}
\noindent\textbf{Stable ID:} 1894\par
\noindent\textbf{Authors:} Dong Wang; David Waugh\par
\noindent\textbf{Original work:} The Transition Probability of the q-TAZRP (q-Bosons) with Inhomogeneous Jump Rates\par
\noindent\textbf{Version:} Official SIGMA 12 (2016) article 037, DOI 10.3842/SIGMA.2016.037; journal PDF and native TeX acquired 2026-10-01, exact bytes pinned by SHA256\par
\noindent\textbf{Selected task scope:} Author1.1: q in(0,1), finite n>=\allowbreak{}1, site conductances a\_i>0,b\_i=\allowbreak{}(1-q)a\_i, ordered X,Y in W\_n on Z and t>=\allowbreak{}0. The specified q-TAZRP stack-top jump rate is a\_i(1-q\textasciicircum{}j). Exact transition formula1.6 is W(X)\textasciicircum{}-1 times the n! sum of n-fold integrals1.5, with inversion amplitudes, extended-product convention and contours enclosing their finite b\_k and q*\allowbreak{}w\_i poles while excluding q\textasciicircum{}-1*\allowbreak{}w\_l poles. No asymptotic limit, invariant-measure normalization or completeness of a new Plancherel transform is selected.\par
\noindent\textbf{Difficulty:} H2: The proof couples the inhomogeneous free evolution with stack boundary conditions through inversion-paired Bethe amplitudes. It then proves the exact delta initial condition by a separate permutation induction and two contour deformations, recovering the q-factorial stack normalization. This completeness argument is substantive integrable probability beyond checking an eigenfunction or importing the homogeneous kernel.\par
\noindent\textbf{Editorial boundary note (not author text):} The complete original main transition-probability Theorem 1.1 and all author general-n proof are retained, with the exact inhomogeneous stack-top rates, ordered chamber/stack weight, inversion amplitudes, extended product and nested-pole contour conventions. All stack boundary/master-equation reduction, inversion-pair cancellation and the full permutation/residue/large-circle delta-initial-condition induction are supplied. The one-particle initial contour and two-particle setup remain uncounted context; no omitted backward-extension converse is assumed. The q-factorial stack normalization is proved, rather than importing Bethe spectral completeness. No step-initial simplification, asymptotic limit, invariant-measure normalization or separate Plancherel claim is counted.\par
\noindent\textbf{Source:} \url{https://sigma-journal.com/2016/037/}\quad DOI: \url{https://doi.org/10.3842/SIGMA.2016.037}\par
\noindent\textbf{License and attribution:} Copyright retained by the named authors. Published by SIGMA under CC BY 4.0: \url{https://creativecommons.org/licenses/by/4.0/}. Source: \url{https://sigma-journal.com/about.html}. Adaptation consists of standalone excerpt formatting, cover and numbering restoration; original author language and mathematical text are retained.\par
\medskip\hrule\medskip\noindent\textbf{Author text begins below.}\par
\setcounter{section}{1}
\setcounter{subsection}{0}
\setcounter{subsubsection}{0}
\setcounter{Theorem}{0}
\setcounter{equation}{0}
% Original source lines 98--190
\subsection[Description of $q$-TAZRP]{Description of $\boldsymbol{q}$-TAZRP}

Throughout this paper, $q \in (0, 1)$ is a constant. Let $n$ particles be on the integer lattice $\intZ$, and denote the position of the $k$-th particle by $x_k$, or $x_k(t)$ if we specify the time $t \geq 0$. We also use $x_k$ to denote the $k$-th particle itself by abuse of notation. We assume that initially $(x_1(0), x_2(0), \dotsc, x_n(0)) =: X(0) \in \Weylchamber^n$, where $\Weylchamber^n = \{ (i_1, \dotsc, i_n) \in \intZ^n \,|\, i_1 \geq i_2 \geq \dotsb \geq i_n \}$. From the dynamics given below, we know that for all $t > 0$ the order $x_1(t) \geq x_2(t) \geq \dotsb \geq x_n(t)$ is kept, or equivalently, $X(t) \in \Weylchamber^n$.

Each particle is at an integer site at any given time, and several consecutive particles can share a site. Suppose $x_{k - 1}(t) > x_k(t) = x_{k + 1}(t) = \dotsb = x_{k + j - 1}(t) = m > x_{k + j}(t)$, then we say that at time $t$, the particles $x_k, x_{k + 1}, \dotsc, x_{k + j - 1}$ are in a stack at $m$ with height $j$, such that $x_k$ is at the top of the stack, $x_{k + i}$ is below $x_{k + i - 1}$ ($i = 1, \dotsc, j - 1$), and $x_{k + j - 1}$ is at the bottom. See Fig.~\ref{fig:particles}.
\begin{figure}[htb]
 \begin{minipage}[t]{0.46\linewidth}
 \centering
 \includegraphics{particles}
 \caption{Particles in $q$-TAZRP. The $k$-th particle is on the top of a stack at $i - 1$ with height~$3$.}
 \label{fig:particles}
 \end{minipage}
 \hfill
 \begin{minipage}[t]{0.48\linewidth}
 \centering
 \includegraphics{particles_jumped}
 \caption{The $k$-th particle jumps to the bottom of the stack at $i$. Until the $k$-th particle jumps, the $(k + 1)$-th and $(k + 2)$-th particles stay put.}
 \label{fig:particles_jumped}
 \end{minipage}
\end{figure}

On each site $i$, we assign a \emph{conductance} $a_i > 0$. Throughout this paper, we denote $b_i = a_i(1 - q)$ for notational simplicity. The dynamics of the particles is given as follows.
If a particle is in a stack but is not on the top of the stack, then it cannot move; otherwise if the particle is on the top of a stack at $m$ with height $j$, it moves to site $m + 1$ at rate $a_m(1 - q^j) = b_m(1 + q + \dotsb + q^{j - 1})$.
The movement occurs regardless of whether the right neighbour site is occupied or not (hence the model is named ``zero range''). If the right neighbour site is occupied, then $x_k$ moves to the bottom of the stack there, and the height of that stack increases by $1$. See Fig.~\ref{fig:particles_jumped} for illustration.

The $q$-TAZRP is clearly a Markov process. Let $X = (x_1, x_2, \dotsc, x_n) \in \Weylchamber^n$ denote a state of the model, then $P(X; t)$, the probability that the model is at state $X$ at time $t$, satisf\/ies the master equation
\begin{gather} \label{eq:abstract_master_eq}
 \frac{d}{dt} P(X; t) = H P(X; t) = \sum_{Y \in \Weylchamber^n} P(Y; t) (c(Y, X) - c(X, Y)),
\end{gather}
where $H$ is the generator of the Markov process, and $c(X, Y)$ is the transition rate for the Markov process from state $X$ to state $Y$. To express the master equation more concretely, for $X = (x_1, \dotsc, x_n) \in \Weylchamber^n$ we denote
\begin{gather*}
 X^{k, -} = (x_1, \dotsc, x_{k - 1}, x_k - 1, x_{k + 1}, \dotsc, x_n), \qquad k = 1, \dotsc, n,
\end{gather*}
and denote
\begin{gather} \label{eq:various_notations}
 \begin{split}
 & \ell(X) :=  \text{$\#$ of the distinct values attained by $x_1, \dotsc, x_n$}, \\
 & v_i(X) :=  \text{the $i$-th largest value attained by $x_1, \dotsc, x_n$}, \\
& p_k(X) :=  \text{$\#$ of $x_1, \dotsc, x_n$ whose value is $k$}, \\
 & n_i(X) :=  p_{v_i}(X), \quad \text{and} \quad N_i(X) := n_1(X) + \dotsb + n_i(X).
 \end{split}
\end{gather}
For instance, if
\begin{gather} \label{eq:general_X}
 X = (\underbrace{s_1, \dotsc, s_1}_{k_1}, \underbrace{s_2, \dotsc, s_2}_{k_2}, \dotsc, \underbrace{s_m, \dotsc, s_m}_{k_m}), \qquad s_1 > s_2 > \dotsb > s_m,
\end{gather}
and let $K_i = k_1 + k_2 + \dotsb + k_i$, then $\ell(X) = m$, $v_i(X) = s_i$, $n_i(X) = k_i$, and $N_i(X) = K_i$. For~$X$ expressed in~\eqref{eq:general_X}, the master equation~\eqref{eq:abstract_master_eq} can be expressed concretely as
\begin{gather} \label{eq:master_eq_concrete}
 \frac{d}{dt} P(X; t) = \sum^m_{i = 1} \big(1 - q^{p_{s_i - 1}(X^{K_i, -})}\big) a_{s_i - 1} P\big(X^{K_i, -}; t\big) - \sum^m_{i = 1} \big(1 - q^{k_i}\big) a_{s_i} P(X; t),
\end{gather}
where $p_{s_i - 1}(X^{K_i, -})$ is the number of particles on site $s_i - 1$ for the state $X^{K_i, -}$ (as def\/ined in~\eqref{eq:various_notations}). Here we remark that if $X \in \Weylchamber^n$ and $k \in \{ 1, \dotsc, n \}$, $X^{k, -}$ may not be in $\Weylchamber^n$. But all~$X^{K_i, -}$ on the right-hand side of~\eqref{eq:master_eq_concrete} are in~$\Weylchamber^n$.

\subsection{Main result}

In this paper we compute $P_Y(X; t)$, the probability that the particles are in state $X \!= \! (x_1, \dotsc, x_n)$ at time $t > 0$, given the initial condition that the particles are in state $Y = (y_1, \dotsc, y_n)$ at time~$0$. Equivalently, $P_Y(X; t)$ is the solution to the master equation \eqref{eq:abstract_master_eq} with initial condition $P_Y(X; 0) = 1$ if $X = Y$ and $P_Y(X; t) = 0$ if $X \in \Weylchamber^n {\setminus} \{ Y \}$.

To state the main theorem, we introduce some notations. First, for~$X \in \Weylchamber^n$ in the form of~\eqref{eq:general_X}, denote
\begin{gather*}
 W(X) = \prod^{\ell(X)}_{i = 1} \prod^{n_i(X)}_{j = 1} \frac{1 - q^j}{1 - q} = \prod^m_{i = 1} [k_i]_q!,
\end{gather*}
where we use the $q$-deformed factorial $[k]_q! = (1 - q)(1 - q^2) \dotsb (1 - q^k)/(1 - q)^k$. Next, let $w_1, \dotsc, w_n$ be variables. For each pair of integers $\alpha < \beta$, denote
\begin{gather*}
 S_{(\beta, \alpha)} = -\frac{q w_{\beta} - w_{\alpha}}{q w_{\alpha} - w_{\beta}}.
\end{gather*}
For a permutation $\sigma \in S_n$, we say $(\beta, \alpha)$ is an \emph{inversion} of $\sigma$ if $\alpha < \beta$ and $\sigma^{-1}(\alpha) > \sigma^{-1}(\beta)$. Then let
\begin{gather*}
 A_{\sigma}(w_1, \dotsc, w_n) = \prod_{(\beta, \alpha) \text{ is an inversion of } \sigma} S_{(\beta, \alpha)}.
\end{gather*}
We write $A_{\sigma} = A_{\sigma}(w_1, \dotsc, w_n)$ if there is no possibility of confusion, but later we occasionally use dif\/ferent variables for $A_{\sigma}$. Note that for $\sigma = \id \in S_n$, since the identity permutation has no inversion, we let $A_{\id} = 1$. For $X = (x_1, \dotsc, x_n) \in \intZ^n$, $Y = (y_1, \dotsc, y_n) \in \intZ^n$ and $\sigma \in S_n$, def\/ine
\begin{gather} \label{eq:integral_Lambda}
 \Lambda_Y(X; t; \sigma) = \left( \prod^n_{k = 1} \frac{-1}{b_{x_k}} \right) \dashint_{C_1} dw_1 \dotsi \dashint_{C_n} dw_n A_{\sigma} \prod^n_{j = 1} \left[ \prodprime^{x_j}_{k = y_{\sigma(j)}} \left( \frac{b_k}{b_k - w_{\sigma(j)}} \right) e^{-w_j t} \right],
\end{gather}
where the notation $\dashint_{C_i} dw_i$ is a shorthand for $(2\pi i)^{-1} \oint_{C_i} dw_i$ and the notation $\prod'$ is an extension of the usual $\prod$ notation such that
\begin{gather*}
 \prodprime^n_{k = m} f(k) =
 \begin{cases}
 \displaystyle \prod^n_{k = m} f(k) & \text{if $n \geq m$}, \\
 1 & \text{if $n = m - 1$}, \\
\displaystyle \prod^{m - 1}_{k = n + 1} \frac{1}{f(k)} & \text{if $n < m - 1$},
 \end{cases}
\end{gather*}
for example, %$\prod^{\prime -3}_{\phantom{\prime}k = 0}$
 $\prod\limits^{-3}_{k = 0}\!\!{\vphantom{\prod}}' f(k) = 1/[f(-1) f(-2)]$. Later we use the identity
\begin{gather*}
 \left( \prodprime^x_{k = m + 1} f(k) \right) \left( \prodprime^m_{k = y} f(k) \right) = \prodprime^x_{k = y} f(k).
\end{gather*}
With respect to any $w_j$, the integrand in \eqref{eq:integral_Lambda} has explicit poles in the form of $w_j = b_k$ for some~$k$, and the factor $A_{\sigma}$ may introduce poles in the form of $w_j = q w_i$ and $w_j = q^{-1} w_l$ for some~$i$ and~$l$ such that $i < j$ and $l > j$. We require that poles $b_k$ and $q w_i$ are enclosed by~$C_j$, but~$q^{-1} w_l$ are not. To be concrete, we can take that $C_j = \{ \lvert z \rvert = R \}$ for all~$j$ with a large enough~$R$, and choose the counter-clockwise orientation.
\begin{Theorem} \label{thm:main}
 Let $X, Y \in \Weylchamber^n$. Then
 \begin{gather} \label{eq:main_result}
 P_Y(X; t) = \frac{1}{W(X)} \sum_{\sigma \in S_n} \Lambda_Y(X; t; \sigma).
 \end{gather}
\end{Theorem}
\setcounter{section}{1}
\setcounter{subsection}{2}
\setcounter{subsubsection}{0}
\setcounter{Theorem}{3}
\setcounter{equation}{8}
% Original source lines 214--593
\section{One particle case} \label{sec:1pt}

In the simplest case that there are only one particle $x_1$, the master equation \eqref{eq:master_eq_concrete} becomes (we write $X = (x)$ as $x$)
\begin{gather*} %\label{eq:master_eq_1_particle}
 \frac{d}{dt} P(x; t) = b_{x - 1} P(x - 1; t) - b_{x} P(x; t).
\end{gather*}
Given the initial condition that the particle is at $y \in \intZ$ at time $t = 0$, that is, $P(y; 0) = 1$ and $P(x; 0) = 0$ if $x \neq y$, we can compute $P_y(x; t)$ by elementary probability. Since the particle jumps only to the right, so $P_y(x; t) = 0$ for all $x < y$. The time that the particle f\/irst jumps away from $y$ is in the exponential distribution with rate~$b_y$, and after it jump away, it cannot be back, so $P_y(y; t) = e^{-b_y t}$. Next, $P_y(y + 1; t)$ is the probability that the particle jumps exactly once from time~$0$ to~$t$, and it is
\begin{gather*}
 \int^t_0 b_y e^{-b_y s} e^{-b_{y + 1} (t - s)} ds = b_y (b_{y + 1} - b_y)^{-1} \big(e^{-b_y t} - e^{-b_{y + 1} t}\big), \qquad \text{if $b_{y + 1} \neq b_y$}.
\end{gather*}
Inductively, we use
\begin{gather*}
 P_y(y + k; t) =   \int^t_0 b_y e^{-b_y s_1} \left( \int^t_{s_1} b_{y + 1} e^{-b_{y + 1} s_2} \left( \int^t_{s_2} \dotsb \vphantom{\int^t_{s_{k - 1}}} \right. \right. \\
\hphantom{P_y(y + k; t) =}{} \times \left. \left. \left( \int^t_{s_{k - 1}} b_{y + k - 1} e^{-b_{y + k - 1} s_k} e^{-b_{y + k} (t - s_k)} ds_k \right) \dotsb ds_3 \right) ds_2 \right) ds_1 \\
\hphantom{P_y(y + k; t)}{} =  \int^t_0 b_y e^{-b_y s_1} P_{y + 1}(y + k; t - s_1) ds_1,
\end{gather*}
and derive
\begin{gather*}
 P_y(x; t) =
 \begin{cases}
 0 & \text{if $x < y$}, \\
 e^{-b_x t} & \text{if $x = y$}, \\
 \displaystyle \left( \prod^{x - 1}_{k = y} b_k \right) \sum^x_{k = y} \left( \prod_{y \leq j \leq x,\ j \neq k} \frac{1}{b_j - b_k} \right) e^{-b_k t} & \text{if $x > y$}.
 \end{cases}
\end{gather*}
Here and later, we assume $b_k$ are distinct, and understand the formulas by \lHopital's rule if some of them are identical. A more convenient way to express $P_y(x; t)$ is ($w = 1 - z^{-1}$)
\begin{subequations}
 \begin{gather}
 P_y(x; t) =   \frac{-1}{b_x} \dashint_C \prodprime^x_{k = y} \frac{b_k}{b_k - w} e^{-wt} dw \label{eq:1_particle_a} \\
\hphantom{P_y(x; t)}{}  =   \frac{1}{b_x} \dashint_{\tilde{C}} \left( \prodprime^x_{k = y} \frac{b_k z}{1 - z(1 - b_k)} \right) e^{(z^{-1} - 1)t} \frac{dz}{z^2}, \label{eq:1_particle_b}
 \end{gather}
\end{subequations}
where the contour $C$ encloses all the poles $b_k$ in positive orientation, and the contour $\tilde{C}$ encloses the pole $0$ in positive orientation. If some $b_k \neq 1$, then $\tilde{C}$ should not enclose poles $(1 - b_k)^{-1}$. Thus the $n = 1$ case of Theorem~\ref{thm:main} is proved.

\begin{Remark}
 Although the expression~\eqref{eq:1_particle_b} of $P_y(x; t)$ seems unnatural, it reduces to the $n = 1$ case \cite[Theorem~2.1]{Korhonen-Lee14} when $b_k = 1$ for all $k \in \intZ$.
\end{Remark}

\section{Two particle case} \label{sec:2_particle_case}

In the $2$-particle case, where the two particles are at $x_1 \geq x_2$, the master equation \eqref{eq:master_eq_concrete} can be written as
\begin{gather} \label{eq:master_eq_2_particles}
 \frac{d}{dt} P((x_1, x_2); t) =
 \begin{cases}
 b_{x_1 - 1} P((x_1 - 1, x_2); t) - b_{x_1} P((x_1, x_2); t) & \\
{} + b_{x_2 - 1} P((x_1, x_2 - 1); t) - b_{x_2} P((x_1, x_2); t),
 & \text{if $x_1 > x_2 + 1$}, \\
 (1 + q) b_xP((x, x); t) - b_{x + 1} P((x + 1, x); t) & \text{if $x_2 = x$} \\
{} + b_{x - 1} P((x + 1, x - 1); t) - b_x P((x + 1, x); t),
& \text{and $x_1 = x + 1$}, \\
 b_{x - 1} P((x, x - 1); t) - (1 + q) b_x P((x, x); t) & \text{if $x_1 = x_2 = x$}.
 \end{cases}\!\!\!\!
\end{gather}
and we want to solve it under the initial condition ($y_1 \geq y_2$)
\begin{gather} \label{eq:init_C_2_particles}
 P((y_1, y_2); 0) = 1, \qquad \text{and} \qquad P((x_1, x_2); 0) = 0 \qquad \text{for all $(x_1, x_2) \in \Weylchamber^2 {\setminus} \{ (y_1, y_2) \}$}.\!\!\!
\end{gather}
To simplify the master equation, we consider functions $u^0((x_1, x_2); t)$ that are dif\/ferentiable in $t \geq 0$, for all $(x_1, x_2) \in \intZ^2$. By contrast, the probabilities $P((x_1, x_2); t)$ can also be viewed as a~dif\/ferentiable function in $t$, but it is only def\/ined when $(x_1, x_2) \in \Weylchamber^2$.

If the functions $u^0((x_1, x_2); t)$ satisfy the dif\/ferential equations (so-called free particle evolution equations)
\begin{gather}
 \frac{d}{dt} u^0((x_1, x_2); t) = b_{x_1 - 1} u^0((x_1 - 1, x_2); t) - b_{x_1} u^0((x_1, x_2); t) \nonumber\\
 \hphantom{\frac{d}{dt} u^0((x_1, x_2); t) =}{} + b_{x_2 - 1} u^0((x_1, x_2 - 1); t) - b_{x_2} u^0((x_1, x_2); t),\label{eq:evolution_eq_2_particles}
\end{gather}
and in addition, they satisfy the relations (so-called boundary condition)
\begin{gather} \label{eq:BC_2_particles}
 b_{x - 1} u^0((x - 1, x); t) = b_{x - 1} q u^0((x, x - 1); t) + b_x (1 - q) u^0((x, x); t) \qquad \text{for all $x \in \intZ$},\!\!\!
\end{gather}
then we have that the functions with $t \geq 0$ and $(x_1, x_2) \in \Weylchamber^2$
\begin{gather} \label{eq:relation_u^0_u_2_particles}
 u((x_1, x_2); t) = \frac{1}{W((x_1, x_2))} u^0((x_1, x_2); t),
 \end{gather}
where
\begin{gather*}
W((x_1, x_2)) =
 \begin{cases}
 1 & \text{if $x_1 > x_2$}, \\
 1 + q & \text{if $x_1 = x_2$},
 \end{cases}
\end{gather*}
satisfy the master equation \eqref{eq:master_eq_2_particles} with $P((x_1, x_2); t) = u((x_1, x_2); t)$. To see it, we note that if $x_1 > x_2 + 1$, then \eqref{eq:master_eq_2_particles} is the same as \eqref{eq:evolution_eq_2_particles} with $P((x_1, x_2); t) = u^0((x_1, x_2); t)$; if $x_1 = x + 1$ and $x_2 = x$, then~\eqref{eq:master_eq_2_particles} is equivalent to~\eqref{eq:evolution_eq_2_particles} with $P((x + 1, x); t) = u^0((x + 1, x); t)$, $P(x + 1, x - 1); t) = u^0((x + 1, x - 1); t)$, and $P((x, x); t) = (1 + q)^{-1} u^0((x_1, x_2); t)$; if $x_1 = x_2 = x$, and we let $P((x, x); t) = u^0((x, ); t)$ and $P((x, x - 1); t) = u^0((x, x - 1); t)$, then the right-hand side of then~\eqref{eq:master_eq_2_particles} is the same as $(1 + q)^{-1}$ times \eqref{eq:evolution_eq_2_particles}, assuming the identity~\eqref{eq:BC_2_particles}. Thus~\eqref{eq:evolution_eq_2_particles},~\eqref{eq:BC_2_particles} and~\eqref{eq:relation_u^0_u_2_particles} imply that $u((x_1, x_2); t)$ satisfy~\eqref{eq:master_eq_2_particles}.

Hence by solving \eqref{eq:evolution_eq_2_particles} with the boundary condition \eqref{eq:BC_2_particles} and the partial initial condition on $\Weylchamber^2$: $u^{(0)}(y_1, y_2; 0) = W((y_1, y_2))$ and $u^{(0)}(x_1, x_2; 0) = 0$ if $(x_1, x_2) \in \Weylchamber^2 {\setminus} \{ (y_1, y_2) \}$, we solve the master equation~\eqref{eq:master_eq_2_particles} with the initial condition~\eqref{eq:init_C_2_particles}.

\begin{Remark} \label{rmk:inverse_2pt}
 It is not obvious that from a solution $u((x_1, x_2); t)$ to \eqref{eq:master_eq_2_particles} where $(x_1, x_2) \in \Weylchamber^2$, we can construct $u^0((x_1, x_2); t)$ for $(x_1, x_2) \in \intZ^2$, that satisfy \eqref{eq:relation_u^0_u_2_particles} if $(x_1, x_2) \in \Weylchamber^2$, as well as solve \eqref{eq:evolution_eq_2_particles} and \eqref{eq:BC_2_particles}. Although it is a digression of our main theorem, for its own interest, we describe the construction is as follows:
\begin{gather*}
 u^0((x_1, x_2); t) = u((x_1, x_2); t) \quad \text{if $x_1 > x_2$}, \qquad u^0((x, x); t) = (1 + q) u((x, x); t), \\
 u^0((x - 1, x); t) = q u((x, x - 1); t) + \frac{b_x}{b_{x - 1}} \big(1 - q^2\big) u((x, x); t),
\end{gather*}
and then inductively with respect to $k \geq 1$, under the assumption that any $u^0((x_1, x_2); t)$ is expressed as linear combinations of $u((z_1, z_2); t)$ if $x_1 \geq x_2 - k$,
\begin{gather}
 u^0((x - (k + 1), x); t) =
 \frac{1}{b_{x - (k + 1)}} \bigg[ \frac{d}{dt} u^0((x - k, x); t) \nonumber\\
 \hphantom{u^0((x - (k + 1), x); t) = }{} - b_{x - 1} u^0((x - k, x - 1); t) + (b_{x - k} + b_x) u^0((x - k, x); t) \bigg].\label{eq:indct_u^0_in_u}
\end{gather}
Note that in \eqref{eq:indct_u^0_in_u}, not only are $u^0((x - k, x - 1); t)$ and $u^0((x - k, x); t)$ expressed by linear combinations of $u((z_1, z_2); t)$, but also $\frac{d}{dt} u^0((x - k, x); t)$ is the linear combinations of $\frac{d}{dt} u((z_1, z_2); t)$, and in turn the linear combinations of $u((z_1, z_2); t)$. For example, for $k = 1$,
\begin{gather*}
   u^0((x - 2, x); t) \\
 =   \frac{1}{b_{x - 2}} \left[ \frac{d}{dt} u^0((x - 1, x); t) - b_{x - 1} u^0((x - 1, x - 1); t) + (b_{x - 1} + b_x) u^0((x - 1, x); t) \right] \\
 =   \frac{1}{b_{x - 2}} \left[ \frac{d}{dt} \left( q u((x, x - 1); t) + \frac{b_x}{b_{x - 1}} \big(1 - q^2\big) u((x, x); t) \right) - b_{x - 1} (1 + q) u((x - 1, x - 1); t) \right. \\
  \hphantom{=\frac{1}{b_{x - 2}}} \left. \vphantom{\frac{d}{dt}} + (b_{x - 1} + b_x) \left( q u((x, x - 1); t) + \frac{b_x}{b_{x - 1}} \big(1 - q^2\big) u((x, x); t) \right) \right] \\
 =   \frac{1 - q^2}{b_{x - 2}} \left[ \left( b_x + q \frac{b^2_x}{b_{x - 1}} \right) u((x, x); t) + b_x u((x, x - 1); t) - b_{x - 1} u((x - 1, x - 1); t) \right. \\
  \hphantom{=\frac{1 - q^2}{b_{x - 2}}} + \left. \frac{q}{1 - q^2} b_{x - 2} u((x, x - 2); t) \right].
\end{gather*}
\end{Remark}

We note that for any meromorphic function $f(w_1, w_2)$, the integral with parameters~$x_1$,~$x_2$,~$t$
\begin{gather*}
 I_f((x_1, x_2); t) = \frac{1}{b_{x_1} b_{x_2}} \dashint_{C_1} \! dw_1 \dashint_{C_2} \! dw_2 f(w_1, w_2) \left( \prodprime^{x_1}_{k = 1} \frac{b_k}{b_k - w_1} \right) \! \left( \prodprime^{x_2}_{k = 1} \frac{b_k}{b_k - w_2} \right) e^{-w_1 t} e^{-w_2 t}
\end{gather*}
satisf\/ies the dif\/ferential equations \eqref{eq:evolution_eq_2_particles} for $u^0((x_1, x_2); t)$. To see it, we compute
\begin{alignat*}{3}
 & \frac{d}{dt} I_f((x_1, x_2); t) =   I_g((x_1, x_2); t), \qquad && \text{where} \quad  g(w_1, w_2) =   -(w_1 + w_2) f(w_1, w_2),& \\
& I_f((x_1 - 1, x_2); t) =   I_{f_1}((x_1, x_2); t), \qquad && \text{where} \quad f_1(w_1, w_2) =   \frac{b_{x_1} - w_1}{b_{x_1 - 1}} f(w_1, w_2),& \\
& I_f((x_1, x_2 - 1); t) =   I_{f_2}((x_1, x_2); t), \qquad && \text{where} \quad  f_2(w_1, w_2) =   \frac{b_{x_2} - w_2}{b_{x_2 - 1}} f(w_1, w_2).&
\end{alignat*}
We then have that for $I_f((x_1, x_2); t)$ to satisfy \eqref{eq:evolution_eq_2_particles}, it suf\/f\/ices to show $g = b_{x_1 - 1} f_1 + b_{x_2 - 1} f_2 - (b_{x_1} + b_{x_2}) f$, which is straightforward to verify.

In the same way as above, we have that $I_f((x_1, x_2); t)$ satisf\/ies the boundary condition \eqref{eq:BC_2_particles} for $u^0((x_1, x_2); t)$ if and only if for all $x \in \intZ$
\begin{gather} \label{eq:integral_boundary_C}
 \dashint_{C_1} dw_1 \dashint_{C_2} dw_2 f(w_1, w_2) (qw_2 - w_1) \left( \prodprime^x_{k = 1} \frac{b_k}{b_k - w_1} \right) \left( \prodprime^x_{k = 1} \frac{b_k}{b_k - w_2} \right) e^{-w_1 t} e^{-w_2 t} = 0.\!\!\!
\end{gather}
Then inspired by the constructions in \cite{Korhonen-Lee14}, we take $f(w_1, w_2) = F_{y_1, y_2}(w_1, w_2)$, where $(y_1, y_2) \in \intZ^2$, and
\begin{gather*}
 F_{y_1, y_2}(w_1, w_2) = \left( \prodprime^0_{k = y_1} \frac{b_k}{b_k - w_1} \right) \left( \prodprime^0_{k = y_2} \frac{b_k}{b_k - w_2} \right) \\
 \hphantom{F_{y_1, y_2}(w_1, w_2) =}{}
 - \frac{qw_1 - w_2}{qw_2 - w_1} \left( \prodprime^0_{k = y_2} \frac{b_k}{b_k - w_1} \right) \left( \prodprime^0_{k = y_1} \frac{b_k}{b_k - w_2} \right).
\end{gather*}
It is straightforward to check that this $I_{F_{y_1, y_2}}((x_1, x_2); t)$ satisf\/ies identity~\eqref{eq:integral_boundary_C}.

Now we consider the value of $I_{F_{y_1, y_2}}((x_1, x_2); 0)$ for $(x_1, x_2)$ and $(y_1, y_2) \in \Weylchamber^2$. First, in the double integral
\begin{gather} \label{eq:int_1_init_2pt}
 \dashint_{C_1} dw_1 \dashint_{C_2} dw_2 \left( \prodprime^{x_1}_{k = y_1} \frac{b_k}{b_k - w_1} \right) \left( \prodprime^{x_2}_{k = y_2} \frac{b_k}{b_k - w_2} \right),
\end{gather}
if $x_2 < y_2$, then the inner integral with respect to $w_2$ vanishes, since the integrand is holomorphic with respect to $w_2$; if $x_2 > y_2$, then by enlarging the contour $C_2$ to $\{ \lvert w_2 \rvert = M \}$ with $M \to \infty$, we still have that the inner integral vanishes. Next, in the double integral
\begin{gather} \label{eq:int_2_init_2pt}
 \dashint_{C_1} dw_1 \dashint_{C_2} dw_2 \frac{qw_1 - w_2}{qw_2 - w_1} \left( \prodprime^{x_1}_{k = y_2} \frac{b_k}{b_k - w_1} \right) \left( \prodprime^{x_1}_{k = y_1} \frac{b_k}{b_k - w_2} \right),
\end{gather}
if $x_2 < y_2$, then $x_2 < y_1$, and we have that the inner integral with respect to $w_2$ vanishes if $x_2 < y_1$, since the integrand has no pole within $C_2$ with respect to $w_2$; if $x_2 > y_2$, then $x_1 > y_2$, and by enlarging the contour $C_1$ to $\{ \lvert w_1 \rvert = M \}$ with $M \to \infty$, we have that the integral on~$C_1$ with respect to $w_1$ vanishes for $x_1 > y_2$. So the integral~\eqref{eq:int_1_init_2pt} vanishes unless $x_2 = x_1$, and the integral~\eqref{eq:int_2_init_2pt} vanishes unless $x_1 = x_2 = y_1 = y_2 = y$. In the case that $x_2 = y_2$, the integral~\eqref{eq:int_1_init_2pt} is equal to $b_{y_1} b_{y_2}$. Similarly, if $x_1 = x_2 = y_1 = y_2$, the integral~\eqref{eq:int_2_init_2pt} is equal to~$q b^2_y$. Since $I_{F_{y_1, y_2}}(x_1, x_2); 0)$ is $(b_{x_1} b_{x_2})^{-1}$ times the dif\/ference between the two integrals in~\eqref{eq:int_1_init_2pt} and~\eqref{eq:int_2_init_2pt}, we have that
\begin{gather*}
 I_{F_{y_1, y_2}}((x_1, x_2); 0) =
 \begin{cases}
 W((x_1, x_2)) & \text{if $(x_1, x_2) = (y_1, y_2)$}, \\
 0 & \text{if $(x_1, x_2) \in \Weylchamber^2 {\setminus} \{ (y_1, y_2) \}$}.
 \end{cases}
\end{gather*}

We conclude that $u^0((x_1, x_2); t) = I_{F_{y_1, y_2}}((x_1, x_2); t)$ is a solution to \eqref{eq:evolution_eq_2_particles} and \eqref{eq:BC_2_particles}, and by~\eqref{eq:relation_u^0_u_2_particles} it yields $u((x_1, x_2); t)$ that is a solution to the master equation \eqref{eq:master_eq_2_particles} with the initial condition \eqref{eq:init_C_2_particles}. Note that our solution $u((x_1, x_2); t)$ is expressed by a double contour integral formula, and it agrees with the $n = 2$ case of~\eqref{eq:main_result} after a simple change of variables. Thus the $n = 2$ case of Theorem~\ref{thm:main} is proved.

\section[$n$ particle case]{$\boldsymbol{n}$ particle case} \label{sec:n_particle_case}

Now we prove Theorem \ref{thm:main} in the general form. Our strategy is the same as in the proof in the $n = 2$ case in Section~\ref{sec:2_particle_case}.

First we consider functions $u^0(X; t)$, where $X = (x_1, \dotsc, x_n) \in \intZ^n$, such that they are dif\/ferentiable functions in $t$, and satisfy the free particle evolution equation
\begin{gather} \label{eq:general_evolution_u^0}
 \frac{d}{dt} u^0(X; t) = \sum^n_{k = 1} \big( b_{x_k - 1} u^0\big(X^{k, -}; t\big) - b_{x_k} u^0(X; t) \big).
\end{gather}
We also consider the boundary conditions that generalize \eqref{eq:BC_2_particles}, such that for all $k = 1, \dotsc, n - 1$
\begin{gather}
 b_{x - 1} u^0((\dotsc, x_{k - 1}, \underbrace{x - 1, x}_{\text{$x_k$ and $x_{k + 1}$}}, x_{k + 2}, \dotsc); t) = q b_{x - 1} u^0((\dotsc, x_{k - 1}, x, x - 1, x_{k + 2}, \dotsc); t) \nonumber\\
 \qquad{}
+ (1 - q) b_x u^0((\dotsc, x_{k - 1}, x, x, x_{k + 2}, \dotsc); t).\label{eq:general_BC}
\end{gather}
Then let for all $X \in \Weylchamber^n$
\begin{gather} \label{eq:general_u_by_u^0}
 u(X; t) = \frac{1}{W(X)} u^0(X; t).
\end{gather}
We have the following result:
\begin{Lemma} \label{lem:BA_coefficients}
 Suppose $u^0(X; t)$ is a solution to the free particle evolution equation~\eqref{eq:general_evolution_u^0} with boundary condition~\eqref{eq:general_BC}. The functions $u(X; t)$ defined by~\eqref{eq:general_u_by_u^0} satisfy the master equation~\eqref{eq:master_eq_concrete} with $P(X; t) = u(X; t)$.
\end{Lemma}

This lemma is essentially equivalent to the (A) $\Leftrightarrow$ (B) part of \cite[Proposition 2.7]{Borodin-Corwin-Sasamoto14}. For the sake of readability, we give the proof below. We also remark that given functions $u(X; t)$ for $X \in \Weylchamber^n$ such that they satisfy \eqref{eq:master_eq_concrete} with $P(X; t) = u(X; t)$, then we can construct $u^0(X; t)$ inductively for all $X \in \intZ^n$, such that~$u^0(X; t)$ satisfy \eqref{eq:general_evolution_u^0} for all~$X \in \intZ^n$, and are related to~$u(X; t)$ by~\eqref{eq:general_u_by_u^0} if $X \in \Weylchamber^n$. The $n = 2$ case is discussed in detail in Remark~\ref{rmk:inverse_2pt}, and we do not give the proof for the general~$n$ case since it is not related to our main result.

\begin{proof}[Proof of Lemma \ref{lem:BA_coefficients}]
 Functions $u(X; t)$ satisfy \eqref{eq:master_eq_concrete} if and only if for all $X$ expressed in~\eqref{eq:general_X},
 \begin{gather} \label{eq:master_eq_equiv}
 \frac{d}{dt} u^0(X; t) = \sum^m_{i = 1} D_i,
 \end{gather}
 where
 \begin{gather*}
 D_i = \frac{W(X)}{W(X^{K_i, -})} \big(1 - q^{p_{s_i - 1}(X^{K_i, -})}\big) a_{s_i - 1} u^0\big(X^{K_i, -}; t\big) - \big(1 - q^{k_i}\big) a_{s_i} u^0(X; t).
 \end{gather*}
 Below we show that for each $i$,
 \begin{gather} \label{eq:check_Bethe_by_i}
 D_i = \sum^{K_i}_{j = K_{i - 1} + 1} \big( b_{s_i - 1} u^0\big(X^{ j, -}; t\big) - b_{s_i} u^0(X; t) \big).
 \end{gather}
 Plugging \eqref{eq:check_Bethe_by_i} into \eqref{eq:master_eq_equiv}, we have that \eqref{eq:master_eq_equiv} is equivalent to \eqref{eq:general_evolution_u^0}, and then prove the lemma.

 In the case $i = m$ or $i < m$ and $s_{i} - 1 > s_{i + 1}$, we have $p_{s_i - 1}(X^{K_i, -}) = 1$ and $W(X)/W(X^{K_i, -})$ $= (1 - q^{k_i})/(1 - q)$. On the other hand, if $i < m$ and $s_i - 1 = s_{i + 1}$, we have $p_{s_i - 1}(X^{K_i, -}) = k_{i + 1} + 1$ and $W(X)/W(X^{K_i, -}) = (1 - q^{k_i})/(1 - q^{k_{i - 1} + 1})$. Then in both cases, we have
 \begin{gather}
 D_i = (1 - q^{k_i}) \big( a_{s_i - 1} u^0\big(X^{K_i, -}; t\big) - a_{s_i} u^0(X; t) \big) \nonumber\\
 \hphantom{D_i}{} = \frac{1 - q^{k_i}}{1 - q} \big( b_{s_i - 1} u^0\big(X^{K_i, -}; t\big) - b_{s_i} u^0(X; t) \big).\label{eq:simple_comp_i}
 \end{gather}
 If $k_i = 1$, then \eqref{eq:simple_comp_i} gives \eqref{eq:check_Bethe_by_i} directly. Otherwise, we use \eqref{eq:general_BC} recursively with $j = K_i - 1$, $K_i - 2, \dotsc, K_{i - 1} + 2, K_{i - 1} + 1$ to derive
 \begin{gather*}
 b_{s_i - 1} u^0\big(X^{j, -}; t\big) = q^{K_i - j} b_{s_i - 1} u^0\big(X^{K_i, -}; t\big) + \big(1 - q^{K_i - j}\big) b_{s_i} u^0(X; t),
 \end{gather*}
 and then verify \eqref{eq:check_Bethe_by_i}.
\end{proof}

Next, we use the idea of the Bethe ansatz to construct the solution to equations~\eqref{eq:general_evolution_u^0} and~\eqref{eq:general_BC} with initial condition
\begin{gather} \label{eq:init_cond_u^0_general}
 u^0(X; 0) =
 \begin{cases}
 W(Y) & \text{if $X = Y \in \Weylchamber^n$}, \\
 0 & \text{if $X \in \Weylchamber^n {\setminus} \{ Y \}$}.
 \end{cases}
\end{gather}
The existence of the solution is not obvious, and it depends on the completeness of the Bethe ansatz for this model, which is so far only proved in~\cite{Borodin-Corwin-Petrov-Sasamoto15} in the homogeneous rate parameter case. Nevertheless, we imitate the construction in~\cite{Korhonen-Lee14} and guess the formula. Then we simply check that the formula is correct.

\begin{Lemma} \label{lem:Bethe_ansatz_solution}
 Let $Y \in \Weylchamber^n$. The function
 \begin{gather*}
 u^0_Y(X; t) = \sum_{\sigma \in S_n} \Lambda_Y(X; t; \sigma)
 \end{gather*}
 satisfies equation~\eqref{eq:general_evolution_u^0}, boundary condition~\eqref{eq:general_BC} and initial condition~\eqref{eq:init_cond_u^0_general}.
\end{Lemma}

\begin{proof}
\textit{Free particle evolution equations \eqref{eq:general_evolution_u^0}.}
 It is straightforward to check that for all $\sigma$, $\Lambda_Y(X; t; \sigma)$ satisfy \eqref{eq:general_evolution_u^0}, so $u^0_Y(X; t)$ also satisf\/ies \eqref{eq:general_evolution_u^0}.

\textit{Boundary conditions \eqref{eq:general_BC}.}
 To check that $u^0_Y(X; t)$ satisf\/ies \eqref{eq:general_BC}, we f\/ix $k \in \{ 1, \dotsc, n - 1 \}$, and denote
 \begin{gather*}
 X(a, b) = (x_1, \dotsc, x_{k - 1}, \!\!\underbrace{a, b}_{x_k \text{ and } x_{k + 1}}\!\! , x_{k + 2}, \dotsc, x_n).
 \end{gather*}
 We divide $S_n$ into $n!/2$ pairs, such that permutations $\mu$ and $\nu$ are in a pair if and only if $\mu = \nu \circ (k, k + 1)$ and $\nu = \mu \circ (k, k + 1)$, that is, $\mu(i) = \nu(i)$ if $i \neq k, k + 1$, and $\mu(k) = \nu(k + 1)$, $\mu(k + 1) = \nu(k)$. Below we assume that $\mu$ and $\nu$ are in a pair such that $\mu(k) = \nu(k + 1) = \alpha < \beta = \mu(k + 1) = \nu(k)$. It is not hard to check that
 \begin{gather*}
 A_{\nu} = A_{\mu} S_{(\beta, \alpha)}.
 \end{gather*}
 For notational simplicity, we denote
 \begin{gather*}
 F(w_1, \dotsc, w_n) =  \frac{1}{b_x} A_{\mu} \prod^n_{l = 1} e^{-w_l t} \prod_{\substack{j = 1, \dotsc, n \\ j \neq k, k + 1}} \frac{-1}{b_{x_j}} \prodprime^{x_j}_{l_j = y_{\mu(j)}} \frac{b_{l_j}}{b_{l_j} - w_{\mu(j)}}, \\
 G(w_1, \dotsc, w_n) =  F(w_1, \dotsc, w_n) \prod^x_{i = y_{\alpha}} \prodprime^x_{j = y_{\beta}} \frac{b_i b_j}{(b_i - w_{\alpha})(b_j - w_{\beta})}.
 \end{gather*}
 Then we have
 \begin{gather}
   b_{x - 1} \left( \Lambda_Y(X(x - 1, x); t; \mu) + \Lambda_Y(X(x - 1, x); t; \nu) \right)\nonumber \\
   \qquad{}
 =   \dashint_{C_1} dw_1 \dotsi \dashint_{C_n} dw_n F(w_1, \dotsc, w_n) \left( \prodprime^{x - 1}_{l_k = y_{\alpha}} \frac{b_{l_k}}{b_{l_k} - w_{\alpha}} \prodprime^x_{l_{k + 1} = y_{\beta}} \frac{b_{l_{k + 1}}}{b_{l_{k + 1}} - w_{\beta}} \right. \nonumber\\
\left. \qquad\quad{} - \frac{qw_{\beta} - w_{\alpha}}{qw_{\alpha} - w_{\beta}} \prodprime^{x - 1}_{l_k = y_{\beta}} \frac{b_{l_k}}{b_{l_k} - w_{\beta}} \prodprime^x_{l_{k + 1} = y_{\alpha}} \frac{b_{l_{k + 1}}}{b_{l_{k + 1}} - w_{\alpha}} \right) \nonumber\\
 \qquad {}=   \dashint_{C_1} dw_1 \dotsi \dashint_{C_n} dw_n G(w_1, \dotsc, w_n) \left( \frac{b_x - w_{\alpha}}{b_x} - \frac{qw_{\beta} - w_{\alpha}}{qw_{\alpha} - w_{\beta}} \frac{b_x - w_{\beta}}{b_x} \right).
\label{eq:boundary_integral_x-1_x}
 \end{gather}
 Similarly,
 \begin{gather}
 b_x \left( \Lambda_Y(X(x, x); t; \mu) + \Lambda_Y(X(x, x); t; \nu) \right) \nonumber\\
 \qquad{}{} = \dashint_{C_1} dw_1 \dotsi \dashint_{C_n} dw_n G(w_1, \dotsc, w_n) \left( 1 - \frac{qw_{\beta} - w_{\alpha}}{qw_{\alpha} - w_{\beta}} \right), \label{eq:boundary_integral_x_x}
 \\
 b_{x - 1} \left( \Lambda_Y(X(x, x - 1); t; \mu) + \Lambda_Y(X(x, x - 1); t; \nu) \right) \nonumber\\
 \qquad{} = \dashint_{C_1} dw_1 \dotsi \dashint_{C_n} dw_n G(w_1, \dotsc, w_n)
  \left( \frac{b_x - w_{\beta}}{b_x} - \frac{qw_{\beta} - w_{\alpha}}{qw_{\alpha} - w_{\beta}} \frac{b_x - w_{\alpha}}{b_x} \right).\label{eq:boundary_integral_x_x-1}
 \end{gather}
 With the integral expressions \eqref{eq:boundary_integral_x-1_x}, \eqref{eq:boundary_integral_x_x} and \eqref{eq:boundary_integral_x_x-1}, and the algebraic identity
 \begin{gather*}
 q \left( \frac{b_x - w_{\beta}}{b_x} - \frac{qw_{\beta} - w_{\alpha}}{qw_{\alpha} - w_{\beta}} \frac{b_x - w_{\alpha}}{b_x} \right) + (1 - q) \left( 1 - \frac{qw_{\beta} - w_{\alpha}}{qw_{\alpha} - w_{\beta}} \right) \\
 \qquad{} = \frac{b_x - w_{\alpha}}{b_x} - \frac{qw_{\beta} - w_{\alpha}}{qw_{\alpha} - w_{\beta}} \frac{b_x - w_{\beta}}{b_x},
 \end{gather*}
 it is clear that
 \begin{gather*}
 b_{x - 1} \sum_{\sigma \in \{ \mu, \nu \}} \Lambda_Y(X(x - 1, x); t; \sigma)   \\
 \qquad{} = q b_{x - 1} \sum_{\sigma \in \{ \mu, \nu \}} \Lambda_Y(X(x, x - 1); t; \sigma) + (1 - q) b_x \sum_{\sigma \in \{ \mu, \nu \}} \Lambda_Y(X(x, x); t; \sigma).
 \end{gather*}
 Since the identity holds for all the $n!/2$ pairs in $S_n$, we sum up them together, and conclude that $u^0_Y(X; t)$ satisf\/ies \eqref{eq:general_BC}.

\textit{Initial condition~\eqref{eq:init_cond_u^0_general}.}
 We check the initial condition inductively. The $n = 1$ case can be done by direct computation, and if we assume that~\eqref{eq:init_cond_u^0_general} holds when the particle number is $n - 1$, we show that it also holds when the particle number is~$n$.

 For any $k = 1, \dotsc, n$, we def\/ine two subsets of $S_n$:
 \begin{gather}
 S_n(k) = \{ \sigma \in S_n \,|\, \sigma(1) = k \}, \qquad \text{and} \nonumber\\
  S^{-1}_n(k) = \{ \sigma \in S_n \,|\, \sigma(k) = 1 \} = \big\{ \sigma \,|\, \sigma^{-1} \in S_n(k) \big\}.\label{eq:defn_S_n(k)}
 \end{gather}
 Then we def\/ine two one-to-one correspondences $\varphi_k\colon S_n(k) \to S_{n - 1}$ and $\psi_k\colon S^{-1}_n(k) \to S_{n - 1}$, such that
for all $i = 1, \dotsc, n - 1$
 \begin{gather}
 \varphi_k(\sigma) =   \tau \quad \text{if\/f} \quad \sigma(i + 1) =
 \begin{cases}
 \tau(i) & \text{if $\tau(i) < k$}, \\
 \tau(i) + 1 & \text{if $\tau(i) \geq k$},
 \end{cases} \nonumber\\
 \psi_k(\sigma) =   \lambda \quad \text{if\/f} \quad \sigma^{-1}(i + 1) =
 \begin{cases}
 \lambda^{-1}(i) & \text{if $\lambda^{-1}(i) < k$}, \\
 \lambda^{-1}(i) + 1 & \text{if $\lambda^{-1}(i) \geq k$}.
 \end{cases}
\label{eq:defn_varphi_k_and_psi_k}
 \end{gather}
 As an example, $\sigma = (123) = \left(
 \begin{smallmatrix}
 1 & 2 & 3 \\
 2 & 3 & 1
 \end{smallmatrix}
 \right) \in S_3$ is in $S_3(2)$ and $S_3(3)$, with $\varphi_2(\sigma) = (12) \in S_2$ and $\psi_3(\sigma) = \id \in S_2$.


 We have for any $\sigma \in S_n$,
\begin{gather}
 A_{\sigma}(w_1, \dotsc, w_n) \nonumber\\
\qquad{} =
 \begin{cases}
 \displaystyle \prod^{k - 1}_{j = 1} \frac{q w_k - w_j}{w_k - q w_j} A_{\tau}(w_1, \dotsc, w_{ k - 1}, w_{k + 1}, w_n) & \text{if $\sigma(1) = k$ and $\varphi_k(\sigma) = \tau$}, \\
 \displaystyle \prod^{k - 1}_{j = 1} \frac{q w_{\sigma(j)} - w_1}{w_{\sigma(j)} - q w_1} A_{\lambda}(w_2, \dotsc, w_n) & \text{if $\sigma(k) = 1$ and $\psi_k(\sigma) = \lambda$}.
 \end{cases}\label{eq:A_sigma_in_A_tau_A_lambda}
 \end{gather}

 We write
 \begin{gather*}
 u^0_Y(X; 0) = \sum^n_{k = 1} I^{(k)}_Y(X), \qquad \text{where} \quad I^{(k)}_Y(X) = \sum_{\sigma \in S_n(k)} \Lambda_Y(X; 0; \sigma),
 \end{gather*}
 and compute $I^{(k)}_Y(X)$ for $k = 1, \dotsc, n$, under the assumption that~\eqref{eq:init_cond_u^0_general} holds when $n$ is replaced by $n - 1$.

 In the simplest case $k = 1$, it is straightforward. Note that if $\sigma \in S_n(1)$, by~\eqref{eq:A_sigma_in_A_tau_A_lambda} $A_{\sigma}(w_1, \dotsc, w_n)$ does not depend on $w_1$, and we have ($(x_1)$ is represented by~$x_1$ as in Section~\ref{sec:1pt})
 \begin{gather*}
  \left( (-1)^n \prod^n_{i = 1} b_{x_i} \right) I^{(1)}_Y(X) \\
 =   \dashint_{C_1} dw_1 \prodprime^{x_1}_{l_1 = y_1} \frac{b_{l_1}}{b_{l_1} - w_1} \sum_{\sigma \in S_n(1)} \dashint_{C_2} dw_2 \dotsi \dashint_{C_n} dw_n A_{\sigma}(w_1, \dotsc, w_n) \prod^n_{j = 2} \prodprime^{x_j}_{l_j = y_{\sigma(j)}} \frac{b_{l_j}}{b_{l_j} - w_{\sigma(j)}} \\
 =   \dashint_{C_1} dw_1 \prodprime^{x_1}_{l_1 = y_1} \frac{b_{l_1}}{b_{l_1} - w_1} \sum_{\tau \in S_{n - 1}} \dashint_{C_2} dw_2 \dotsi \dashint_{C_n} dw_n A_{\tau}(w_2, \dotsc, w_n) \prod^{n - 1}_{j = 1} \prodprime^{x_{j + 1}}_{l_j = y_{\tau(j) + 1}} \frac{b_{l_j}}{b_{l_j} - w_{\tau(j) + 1}} \\
 =  \left( (-1)^n \prod^n_{i = 1} b_{x_i} \right) u^0_{y_1}(x_1; 0) u^0_{(y_2, \dotsc, y_n)}((x_2, \dotsc, x_n); 0).
 \end{gather*}
 So by the inductive assumption,
 \begin{gather*}
 I^{(1)}_Y(X) =
 \begin{cases}
 W((x_2, \dotsc, x_n)) & \text{if $X = Y$}, \\
 0 & \text{if $X \in \Weylchamber^n {\setminus} \{ Y \}$}.
 \end{cases}
 \end{gather*}

 In the general case $k > 1$, for all $\sigma \in S_n(k)$, if $S_{n - 1} \ni \tau = \varphi_k(\sigma)$, we have
 \begin{gather}
 \left( (-1)^n \prod^n_{i = 1} b_{x_i} \right) \Lambda_Y(X; 0; \sigma) \nonumber\\
\qquad{} = \dashint_{C_1} dw_1 \dotsi \dashint_{C_{k - 1}} dw_{k - 1} \dashint_{C_{k + 1}} dw_{k + 1} \dotsi \dashint_{C_n} dw_n A_{\tau}(w_1, \dotsc, w_{k - 1}, w_{k + 1}, \dotsc, w_n) \nonumber\\
 \qquad\quad{} \times \prod^n_{j = 2} \prodprime^{x_j}_{l_j = y_{\sigma(j)}} \frac{b_{l_j}}{b_{l_j} - w_{\sigma(j)}} \left( \dashint_{C_k} \prod^{k - 1}_{i = 1} \frac{q w_k - w_j}{q w_j - w_k} \prodprime^{x_1}_{l_1 = y_k} \frac{b_{l_1}}{b_{l_1} - w_k} dw_k \right).\label{eq:Lambda_YX;0;sigma_first}
 \end{gather}
 The inner integral vanishes if $x_1 > y_k$, so $\Lambda_Y(X; 0; \sigma) = 0$ if $x_1 > y_k$. To see it, we deform the contour $C_k$ into a circle centred at $0$ with radius $R_k \to \infty$ while keeping the other $C_j$ ($j \neq k$) in place, and use the estimate that the integrand of the inner integral is $\bigO\big(w^{y_k - x_1 - 1}_k\big)$ as $w_k \to \infty$. Note that this deformation does not change the integral on the right-hand side of~\eqref{eq:Lambda_YX;0;sigma_first}, since the only requirement for $C_k$ is that it is big enough to enclose all possible poles~$b_{l_1}$.

 Next, we assume that $\sigma \in S^{-1}_n(m)$, and $\psi_m(\sigma) = \lambda$. Then we have
 \begin{gather}
 \left( (-1)^n \prod^n_{i = 1} b_{x_i} \right) \Lambda_Y(X; 0; \sigma) = \dashint_{C_2} dw_2 \dotsi \dashint_{C_n} dw_n A_{\lambda}(w_2, \dotsc, w_n) \prod^n_{j = 2} \prodprime^{x_{\sigma^{-1}(j)}}_{l_j = y_j} \frac{b_{l_j}}{b_{l_j} - w_j}\nonumber \\
 \hphantom{\left( (-1)^n \prod^n_{i = 1} b_{x_i} \right) \Lambda_Y(X; 0; \sigma) =}{}
 \times \left( \dashint_{C_1} \prod^{k - 1}_{j = 1} \frac{q w_{\sigma(j)} - w_1}{w_{\sigma(j)} - q w_1} \prodprime^{x_m}_{l_1 = y_1} \frac{b_{l_1}}{b_{l_1} - w_1} dw_1 \right).\label{eq:Lambda_YX;0;sigma_second}
 \end{gather}
 If $x_m < y_1$, then the inner integral in~\eqref{eq:Lambda_YX;0;sigma_second} vanishes, for its integrand has no pole within~$C_1$. So~$\Lambda_Y(X; 0; \sigma) = 0$ if~$x_m < y_1$.

 Since $x_1 \geq \dotsb \geq x_m$ and $y_1 \geq \dotsb \geq y_k$, by argument above, we see that $\Lambda_Y(X; 0; \sigma) = 0$ unless $x_1 = \dotsb = x_m = y_1 = \dotsb = y_k$, and in particular, $x_1 = y_1 = \dotsb = y_k$.

 When $x_1 = y_1 = \dotsb = y_k$, by deforming the contour $C_k$ into a big circle centred at $0$ with radius $R_k \to \infty$, the inner integral in \eqref{eq:Lambda_YX;0;sigma_first} can be explicitly computed as
 \begin{gather} \label{eq:inner_integral}
 \dashint_{C_k} \prod^{k - 1}_{i = 1} \frac{q w_k - w_j}{w_k - q w_j} \frac{b_{x_1}}{b_{x_1} - w_k} dw_k = -q^{k - 1} b_{x_1},
 \end{gather}
 since its integrand is $-q^{k - 1} b_{x_1} w^{-1}_k \big(1 + \bigO\big(w^{-1}_k\big)\big)$. Plugging in~\eqref{eq:inner_integral} into~\eqref{eq:Lambda_YX;0;sigma_first}, we have
 \begin{gather}
   \left( (-1)^n \prod^n_{i = 1} b_{x_i} \right) \Lambda_Y(X; 0; \sigma)
 =  \dashint_{C_1} dw_1 \dotsi \dashint_{C_{k - 1}} dw_{k - 1} \dashint_{C_{k + 1}} dw_{k + 1} \dotsi \nonumber\\
 \qquad\quad{}\times \dashint_{C_n} dw_n A_{\tau}(w_1, \dotsc, w_{k - 1}, w_{k + 1}, \dotsc, w_n)  \prod^n_{j = 2} \prodprime^{x_j}_{l_j = y_{\sigma(j)}} \frac{b_{l_j}}{b_{l_j} - w_{\sigma(j)}} q^{k - 1} b_{x_1} \nonumber\\
\qquad{}
 =  q^{k - 1} \left( (-1)^n \prod^n_{i = 1} b_{x_i} \right) \Lambda_{(y_2, \dotsc, y_n)}((x_2, \dotsc, x_n); 0; \tau),\label{eq:evanluation_x_1=y_k}
  \end{gather}
 where $\tau = \varphi_k(\sigma)$. Thus summing up with $\sigma \in S_n(k)$, we have that $\tau$ runs over $S_{n - 1}$ and
 \begin{gather*} %\label{eq:evanluation_x_1=y_k_I^k}
 I^{(k)}_Y(X) = q^{k - 1} u^0_{(y_2, \dotsc, y_n)}((x_2, \dotsc, x_n); 0).
 \end{gather*}
 If $X \neq Y$, then the assumption $x_1 = y_1 = \dotsb = y_k$ implies $(x_2, \dotsc, x_n) \neq (y_2, \dotsc, y_n)$, and by inductive assumption have that $I^{(k)}_Y(X) = q^{k - 1} \cdot 0 = 0$. Otherwise, $X = Y$, and then the assumption $x_1 = y_1 = \dotsb = y_k$ holds only if $k \leq k_1$, if $X$ is expressed in the form of~\eqref{eq:general_X}.

 The result derived above can be summarized as follows: If $X$ is expressed by~\eqref{eq:general_X}, then by the inductive assumption,
 \begin{gather*}
 I^{(k)}_X(X) =
 \begin{cases}
 0 & \text{if $X \in \Weylchamber^n {\setminus} \{ Y \}$}, \\
 0 & \text{if $X = Y$ and $k > k_1$}, \\
 q^{k - 1} W((x_2, \dotsc, x_n)) & \text{if $X = Y$ and $k \leq k_1$},
 \end{cases}
 \end{gather*}
 and hence $u^0_Y(X) = 0$ if $X \in \Weylchamber^n {\setminus} \{ Y \}$ and $u^0_X(X) = \big(1 + q + \dotsb + q^{k_1 - 1}\big) W((x_2, \dotsc, x_n)) = W(X)$, and we prove~\eqref{eq:init_cond_u^0_general}.
\end{proof}

Combining Lemmas \ref{lem:Bethe_ansatz_solution} and \ref{lem:BA_coefficients}, we prove Theorem \ref{thm:main} in the general $n$ particle case.
\begin{thebibliography}{99}
\bibitem[4]{Borodin-Corwin-Petrov-Sasamoto15}
Borodin A., Corwin I., Petrov L., Sasamoto T., Spectral theory for the
 {$q$}-{B}oson particle system, \href{http://dx.doi.org/10.1112/S0010437X14007532}{\textit{Compos. Math.}} \textbf{151} (2015),
 1--67, \href{http://arxiv.org/abs/1308.3475}{arXiv:1308.3475}.

\bibitem[5]{Borodin-Corwin-Sasamoto14}
Borodin A., Corwin I., Sasamoto T., From duality to determinants for
 {$q$}-{TASEP} and {ASEP}, \href{http://dx.doi.org/10.1214/13-AOP868}{\textit{Ann. Probab.}} \textbf{42} (2014),
 2314--2382, \href{http://arxiv.org/abs/1207.5035}{arXiv:1207.5035}.

\bibitem[9]{Korhonen-Lee14}
Korhonen M., Lee E., The transition probability and the probability for the
 left-most particle's position of the {$q$}-totally asymmetric zero range
 process, \href{http://dx.doi.org/10.1063/1.4851758}{\textit{J.~Math. Phys.}} \textbf{55} (2014), 013301, 15~pages,
 \href{http://arxiv.org/abs/1308.4769}{arXiv:1308.4769}.

\end{thebibliography}
\end{document}
